\newpage
\appendix
\onecolumn

\section{Comparison of Statistics Between Point Estimation and Discounted Beta--Bernoulli Estimation}
\FloatBarrier

We present Table~\ref{tab:estimator_comparison} that summarizes the statistics of point estimation and Discounted Beta--Bernoulli estimation, as discussed in Sections~\ref{subsec:problem_formulation} and~\ref{subsec:grpohistbeta}, in a concise and accessible manner.


\begin{table}[htbp]
\centering
\renewcommand{\arraystretch}{2.0}   
\setlength{\tabcolsep}{4pt}         

\caption{Comparison between point estimation and DBB estimation for reward distribution modeling. We denote the Bernoulli distribution by $\mathrm{Bern}(\cdot)$.}
\label{tab:estimator_comparison}
\begin{tabular}{l|c|c}
\toprule
 & {\small\textbf{Point Estimation}} & {\small\textbf{DBB Estimation}} \\
\midrule

$X_{\tau,i}$
& $X_{\tau,i} \sim \mathrm{Bern}(p_\tau)$
& $X_{\tau,i} \sim \mathrm{Bern}(p_\tau)$ \\


$p_\tau$
& fixed but unknown
& $p_\tau \sim \mathrm{Beta}(\alpha_\tau,\beta_\tau)$ \\

$\hat{p}_\tau$
& $\displaystyle \hat{p}_\tau^{\mathrm{pt}} = \frac{1}{N}\sum_{i=1}^N X_{\tau,i}$
& $\displaystyle \hat{p}_\tau^{\mathrm{dbb}} = \frac{\alpha_\tau}{\alpha_\tau+\beta_\tau}$ \\

$\widehat{\mathrm{Var}}(X_\tau)$
& $\displaystyle \frac{N}{N-1}\,\hat{p}_\tau^{\mathrm{pt}}\!\left(1-\hat{p}_\tau^{\mathrm{pt}}\right)$
& $\displaystyle \hat{p}_\tau^{\mathrm{dbb}}\!\left(1-\hat{p}_\tau^{\mathrm{dbb}}\right)$ \\


$\mathbb{E}[\hat{p}_\tau]$
& $p_\tau$
& $w\mu_{\tau-1} + (1-w)p_\tau$ \\

$\mathrm{Bias}(\hat{p}_\tau)$
& $0$
& $w(\mu_{\tau-1} - p_\tau)$ \\

$\mathrm{Var}(\hat{p}_\tau)$
& $\displaystyle \frac{p_\tau(1-p_\tau)}{N}$
& $\displaystyle (1-w)^2\frac{p_\tau(1-p_\tau)}{N}$ \\

\bottomrule
\end{tabular}
\end{table}

\FloatBarrier



\section{Derivation of the Mean, Variance, and MSE of the DBB Estimator}
\label{app:hb-derivation}

In this section, we provide a detailed derivation of Equations~\eqref{eq:hb_expectation_pseq} and~\eqref{eq:hb_variance_pseq} in the main text.
To avoid ambiguity caused by implicit conditioning, we re-express the analysis in terms of the conditional distribution given the sequence of true reward probabilities $\{p_1, p_2, \dots, p_\tau\}$.

\subsection{Expansion of Posterior Parameters under Historical Rewards}

At each training step $k$, rollout rewards satisfy
\begin{equation}
S_k \mid p_k \sim \mathrm{Binomial}(N, p_k),
\qquad
X_{k,i} \mid p_k \sim \mathrm{Bernoulli}(p_k),
\end{equation}
and the DBB updates are given by
\begin{equation}
\alpha_k = \lambda \alpha_{k-1} + S_k,
\qquad
\beta_k = \lambda \beta_{k-1} + (N - S_k).
\end{equation}

Unrolling the recursion yields
\begin{equation}
\alpha_\tau
=
\lambda^\tau \alpha_0
+
\sum_{k=1}^{\tau} \lambda^{\tau-k} S_k,
\end{equation}
and the total mass
\begin{equation}
H_\tau
\;\triangleq\;
\alpha_\tau + \beta_\tau
=
\lambda^\tau(\alpha_0 + \beta_0)
+
N \sum_{k=1}^{\tau} \lambda^{\tau-k}.
\end{equation}

Crucially, when $N$ and $\lambda$ are fixed, $H_\tau$ is deterministic and does not depend on the random variables $\{S_k\}_{k=1}^{\tau}$.
All stochasticity in the estimator arises from $\alpha_\tau$ alone.

\subsection{Mean of the DBB Estimator}

The DBB estimator is defined as
\begin{equation}
\hat{p}_\tau^{\mathrm{dbb}} = \frac{\alpha_\tau}{H_\tau}.
\end{equation}
Since $H_\tau$ is deterministic,
\begin{equation}
\mathbb{E}\!\left[\hat{p}_\tau^{\mathrm{dbb}} \mid p_{1:\tau}\right]
=
\frac{\mathbb{E}[\alpha_\tau \mid p_{1:\tau}]}{H_\tau}.
\end{equation}

Using $\mathbb{E}[S_k \mid p_k] = N p_k$, we obtain
\begin{equation}
\mathbb{E}[\alpha_\tau \mid p_{1:\tau}]
=
\lambda^\tau \alpha_0
+
N \sum_{k=1}^{\tau} \lambda^{\tau-k} p_k.
\end{equation}
Therefore,
\begin{equation}
\boxed{
\mathbb{E}\!\left[\hat{p}_\tau^{\mathrm{dbb}} \mid p_{1:\tau}\right]
=
\frac{\lambda^\tau \alpha_0 + N \sum_{k=1}^{\tau} \lambda^{\tau-k} p_k}
{\lambda^\tau(\alpha_0+\beta_0)+N \sum_{k=1}^{\tau} \lambda^{\tau-k}}
}
\label{eq:hb-mean}
\end{equation}

For interpretability, define weights
\begin{equation}
c_0 \triangleq \frac{\lambda^\tau(\alpha_0+\beta_0)}{H_\tau},
\qquad
c_k \triangleq \frac{N\lambda^{\tau-k}}{H_\tau}, \quad k=1,\dots,\tau,
\qquad
p_0 \triangleq \frac{\alpha_0}{\alpha_0+\beta_0}.
\end{equation}
Then the estimator mean can be written as
\begin{equation}
\boxed{
\mathbb{E}\!\left[\hat{p}_\tau^{\mathrm{dbb}} \mid p_{1:\tau}\right]
=
\sum_{k=0}^{\tau} c_k p_k,
\qquad
\sum_{k=0}^{\tau} c_k = 1,
}
\label{eq:hb-mean-weighted}
\end{equation}
revealing an exponentially weighted average of historical reward probabilities.

\subsection{Variance of the DBB Estimator}

Again using the determinism of $H_\tau$,
\begin{equation}
\mathrm{Var}\!\left(\hat{p}_\tau^{\mathrm{dbb}} \mid p_{1:\tau}\right)
=
\frac{\mathrm{Var}(\alpha_\tau \mid p_{1:\tau})}{H_\tau^2}.
\end{equation}

Since $\{S_k\}$ are independent conditioned on $\{p_k\}$ and
$\mathrm{Var}(S_k \mid p_k)=N p_k(1-p_k)$,
\begin{equation}
\mathrm{Var}(\alpha_\tau \mid p_{1:\tau})
=
\sum_{k=1}^{\tau} \lambda^{2(\tau-k)} N p_k(1-p_k).
\end{equation}
Hence,
\begin{equation}
\boxed{
\mathrm{Var}\!\left(\hat{p}_\tau^{\mathrm{dbb}} \mid p_{1:\tau}\right)
=
\frac{\sum_{k=1}^{\tau} \lambda^{2(\tau-k)} N p_k(1-p_k)}
{\left(\lambda^\tau(\alpha_0+\beta_0)+N \sum_{k=1}^{\tau} \lambda^{\tau-k}\right)^2}
}
\label{eq:hb-var}
\end{equation}

\subsection{Bias and Mean Squared Error of the DBB Estimator}

By definition,
\begin{equation}
\mathrm{MSE}\!\left(\hat{p}_\tau^{\mathrm{dbb}} \mid p_{1:\tau}\right)
=
\mathrm{Bias}^2 + \mathrm{Var}.
\end{equation}

Using Eq.~\eqref{eq:hb-mean-weighted}, the conditional bias is
\begin{equation}
\boxed{
\mathrm{Bias}\!\left(\hat{p}_\tau^{\mathrm{dbb}} \mid p_{1:\tau}\right)
=
\sum_{k=0}^{\tau-1} c_k (p_k-p_\tau),
}
\label{eq:hb-bias}
\end{equation}
where the $k=\tau$ term vanishes since $p_\tau-p_\tau=0$.

Combining Eqs.~\eqref{eq:hb-var} and~\eqref{eq:hb-bias}, we obtain
\begin{equation}
\boxed{
\mathrm{MSE}\!\left(\hat{p}_\tau^{\mathrm{dbb}} \mid p_{1:\tau}\right)
=
\left(
\sum_{k=0}^{\tau-1} c_k (p_k-p_\tau)
\right)^2
+
\frac{\sum_{k=1}^{\tau} \lambda^{2(\tau-k)} N p_k(1-p_k)}{H_\tau^2},
}
\label{eq:hb-mse}
\end{equation}
where $H_\tau=\lambda^\tau(\alpha_0+\beta_0)+N \sum_{k=1}^{\tau} \lambda^{\tau-k}$.


\section{Implementation Details}
\label{ap:A}

\FloatBarrier
\begin{table}[htbp]
\centering
\small
\setlength{\tabcolsep}{6pt}
\caption{Common Training hyperparameter settings for all experiments.}
\label{tab:hyperparams}
\begin{tabular}{lcc}
\toprule
\textbf{Hyperparameter} 
& \textbf{Qwen3-1.7B-Base} 
& \textbf{Qwen3-8B-Base} \\

\midrule
\multicolumn{3}{c}{\textbf{Training Configuration}} \\
\midrule
Training batch size        & 128  & 128 \\
Mini-batch size            & 64   & 64  \\
Number of epochs           & 4    & 4 \\
Total gradient steps       & 1080  & 1080 \\
Samples per prompt         & 8    & 8  \\
Max response length        & 4096 & 8192 \\

\midrule
\multicolumn{3}{c}{\textbf{Sampling Configuration}} \\
\midrule
Training temperature       & 1.0  & 1.0 \\
Training top-$p$           & 1.0  & 1.0 \\
Validation temperature     & 0.6  & 0.6 \\
Validation top-$p$         & 0.95  & 0.95 \\

\midrule
\multicolumn{3}{c}{\textbf{Optimization}} \\
\midrule

Optimizer                  & AdamW & AdamW \\
Learning rate              & $1\times10^{-6}$ & $1\times10^{-6}$ \\
LR warmup steps            & 0    & 0 \\
LR scheduler               & constant & constant \\

\bottomrule
\end{tabular}
\end{table}
\FloatBarrier


In this section, we describe the experimental setup details for the RLVR algorithm with the DBB process (RLVR-DBB) and all baseline methods.
All training experiments are conducted using the \texttt{verl} framework, and \texttt{Math-Verify} is employed to extract and normalize final answers from model responses.
Training Qwen3-1.7B-Base is performed using 4$\times$H200 GPUs, while Qwen3-8B-Base is trained using 8$\times$H200 GPUs.
Due to its substantially higher GPU memory requirements, RePO is trained using 8$\times$H200 GPUs for both model scales.

All training experiments are conducted using the DAPO-Math-17k dataset.
For training stability, we filter prompts to a maximum length of 1024 tokens.
In addition, to facilitate analysis, we further restrict the dataset so that its size is an integer multiple of the rollout batch size, which is 128.
As a result, out of the original 17,391 prompts, 17,280 are used for training.


The hyperparameters shared across all training experiments are summarized in Table~\ref{tab:hyperparams}.
Our experimental setup follows configurations that are most commonly adopted in recent RLVR studies.
Following DAPO, we remove the KL regularization term during training.
In addition, under the clip-higher strategy, we use a clipping range of $(\text{clip}_{\text{low}}, \text{clip}_{\text{high}}) = (0.2, 0.28)$ for all methods except RLVR-DBB.

For RLVR-DBB, the importance ratios tend to be larger than those of naive methods.
As a result, using the same clipping range would cause a substantially larger fraction of meaningful update signals to be clipped.
To avoid overly aggressive clipping, we therefore adopt a wider clipping range of $(0.98, 0.98)$ specifically for RLVR-DBB.

For RePO, which introduces additional hyperparameters and configuration options, we follow the optimal settings recommended in the original paper.
Specifically, we set the number of replay samples to 8, use a replay cache size of 16, and adopt the reward-oriented replay strategy.

For both evaluation and validation, we use the sampling parameters specified in Table~\ref{tab:hyperparams}.
To reduce training overhead, the validation set is restricted to AIME24/25, AMC24, Minerva, and MATH500.
Evaluation on OlympiadBench and the out-of-distribution benchmarks (MMLU-Pro, GPQA-Diamond, and Big-Bench Hard) is performed only on the best checkpoint selected for each experiment.



